\documentclass[11pt]{article}

\usepackage[margin=1in]{geometry}
\usepackage{amsmath,amssymb}
\usepackage{array}
\usepackage{enumitem}

\title{Lecture 6 Notes: \\ Formalism: Do Mathematical Statements Mean Anything?}
\author{Philosophy of Mathematics}
\date{}

\begin{document}

\maketitle

\section*{Central Theme}

Formalism is the family of views according to which mathematics is best understood
through symbols, rules, axioms, and formal systems.

\[
\boxed{
\text{Mathematics} = \text{symbols} + \text{rules?}
}
\]

The guiding question:

\[
\boxed{
\text{Is mathematics about anything, or is it just rule-governed symbol manipulation?}
}
\]

Formalism tries to avoid metaphysical and epistemological problems by treating
mathematics as formal activity rather than as knowledge of abstract objects.

\section*{1. The Formalist Impulse}

Much mathematical practice looks like symbol manipulation according to rules.

For example, if we have:

\[
a \times b = c,
\]

then we may write:

\[
b \times a = c.
\]

And if we also have:

\[
a \neq 0,
\]

then we may write:

\[
c/a = b.
\]

The formalist asks whether this is not merely how mathematics is represented, but what
mathematics essentially is.

\subsection*{Teaching Point}

Formalism seizes on the rule-governed, symbolic character of mathematics and asks
whether anything more is needed.

\subsection*{Whiteboard}

\[
\text{Mathematics}
\quad =
\quad
\text{symbols}
+
\text{rules for manipulating symbols}
\]

\[
\text{Do the symbols have to mean anything?}
\]

\section*{2. Why Formalism Is Tempting}

Formalism becomes especially tempting when mathematics introduces entities that seem
difficult to interpret.

Examples:

\[
\text{negative numbers}
\]

\[
\text{irrational numbers}
\]

\[
\text{transcendental numbers}
\]

\[
\text{imaginary numbers}
\]

\[
\text{points at infinity}
\]

Instead of asking what these entities really are, the formalist may say that mathematics
only tells us how to use the corresponding symbols.

\subsection*{Teaching Point}

Formalism offers a way to avoid metaphysical anxiety about mysterious mathematical
objects.

\subsection*{Whiteboard}

\[
\text{What is } i?
\]

\[
\text{Formalist answer:}
\quad
\text{a symbol governed by rules, e.g. } i^2=-1.
\]

\section*{3. Term Formalism}

The first version of formalism is term formalism.

\[
\boxed{
\text{Mathematical objects are symbols or terms.}
}
\]

On this view:

\[
2 = \text{the numeral } ``2"
\]

\[
8+2i = \text{the symbol } ``8+2i"
\]

So mathematics has a subject matter, but its subject matter is linguistic.

\[
\begin{array}{rcl}
\text{What is mathematics about?} &:& \text{symbols} \\
\text{What are numbers?} &:& \text{numerals or symbol-types} \\
\text{How do we know mathematics?} &:& \text{by knowing rules for symbols}
\end{array}
\]

\subsection*{Teaching Point}

Term formalism does not say that mathematics is about nothing. It says that mathematics
is about its own signs.

\subsection*{Whiteboard}

\[
\textbf{Term Formalism}
\]

\[
\text{numbers} = \text{numerals}
\]

\[
\text{mathematical objects} = \text{mathematical terms}
\]

\section*{4. Tokens and Types}

Term formalism immediately faces a problem.

Consider:

\[
0=0.
\]

This cannot mean that one physical mark is identical to another physical mark. The two
marks are distinct pieces of ink, chalk, or pixels.

So the term formalist may shift from tokens to types.

\[
\begin{array}{c|c}
\textbf{Token} & \textbf{Type} \\
\hline
\text{particular physical mark} & \text{repeatable abstract form} \\
\text{this ink mark} & \text{the symbol-form } ``0" \\
\text{created and destroyed} & \text{shared by many tokens}
\end{array}
\]

But types are abstract objects. So term formalism may reintroduce the very kind of
abstract ontology it hoped to avoid.

\subsection*{Teaching Point}

If numbers are symbol-types, then term formalism still owes us an account of abstract
objects.

\subsection*{Whiteboard}

\[
0=0
\]

\[
\text{token identity? No.}
\]

\[
\text{type identity? Maybe.}
\]

\[
\text{But types are abstract.}
\]

\section*{5. Frege's Critique of Term Formalism}

Frege presses the problem with ordinary equations.

Consider:

\[
5+7=6+6.
\]

If mathematical terms denote themselves, then this says:

\[
``5+7" = ``6+6".
\]

But that is false. The symbols are different.

So the term formalist may reinterpret equality as intersubstitutability:

\[
``5+7" \text{ may be substituted for } ``6+6".
\]

But then equality no longer means identity.

\subsection*{Teaching Point}

Term formalism cannot preserve the ordinary meaning of mathematical equations.

\subsection*{Whiteboard}

\[
5+7=6+6
\]

\[
\text{Not: } ``5+7" = ``6+6"
\]

\[
\text{Maybe: } ``5+7" \leftrightarrow ``6+6"
\text{ by rule}
\]

\[
\text{But then } = \text{ has changed meaning.}
\]

\section*{6. Real Numbers and General Theorems}

Term formalism has special trouble with real numbers.

Many real numbers have no finite name. One might try to identify \(\pi\) with its decimal
expansion:

\[
3.14159\ldots
\]

But the decimal expansion is infinite, not an ordinary concrete symbol.

If the term formalist appeals to limits of finite decimals, then the view begins to look like
ordinary real analysis rather than formalism.

Term formalism also struggles with general mathematical theorems:

\[
\text{There are infinitely many primes.}
\]

\[
\text{Every continuous function on a closed interval is bounded.}
\]

\[
\text{Every nonconstant polynomial over } \mathbb{C} \text{ has a root.}
\]

These do not seem to be merely about symbols.

\subsection*{Teaching Point}

Term formalism may handle simple calculations, but it struggles with higher mathematics
and general mathematical propositions.

\subsection*{Whiteboard}

\[
\pi = 3.14159\ldots
\]

\[
\text{Is this a symbol?}
\]

\[
\text{What about theorems about all real numbers?}
\]

\section*{7. Game Formalism}

The second version of formalism is game formalism.

\[
\boxed{
\text{Mathematics is like chess: a game played with symbols according to rules.}
}
\]

On this view, mathematical symbols are either meaningless or their meaning is irrelevant.

\[
\begin{array}{rcl}
\text{What is mathematics about?} &:& \text{nothing} \\
\text{What are numbers?} &:& \text{nothing, or no concern} \\
\text{What is mathematical knowledge?} &:& \text{knowledge of legal moves}
\end{array}
\]

A theorem is like a reachable position in a game.

\[
\text{axioms}
\quad \longrightarrow \quad
\text{legal moves}
\quad \longrightarrow \quad
\text{theorem}
\]

\subsection*{Teaching Point}

Game formalism avoids ontology by denying that mathematics needs a subject matter.

\subsection*{Whiteboard}

\[
\textbf{Game Formalism}
\]

\[
\text{axioms} + \text{rules} = \text{legal derivations}
\]

\[
\text{mathematical theorem} = \text{reachable position}
\]

\section*{8. Frege's Critique of Game Formalism}

Frege's main objection is the problem of applicability.

Chess is rule-governed, but no one expects chess to explain planetary motion, engineering,
electricity, or quantum mechanics.

Mathematics, by contrast, is deeply applicable to nature.

\[
\begin{array}{c|c}
\textbf{Chess} & \textbf{Mathematics} \\
\hline
\text{rule-governed} & \text{rule-governed} \\
\text{no scientific application expected} & \text{scientifically indispensable} \\
\text{no external truth needed} & \text{seems to express truth}
\end{array}
\]

If mathematics is only a meaningless game, why does it work so well?

\subsection*{Teaching Point}

Frege argues that applicability requires meaning. Arithmetic can be applied because its
equations express thoughts.

\subsection*{Whiteboard}

\[
\text{If mathematics is like chess,}
\]

\[
\text{why does mathematics apply to science?}
\]

\[
\boxed{
\text{Meaning seems needed for application.}
}
\]

\section*{9. Instrumentalism Analogy}

Game formalism resembles instrumentalism in philosophy of science.

Instrumentalism says that theories about unobservable entities are instruments for
prediction, not literal descriptions of reality.

Likewise, game formalism says that mathematics is an instrument made of formal
symbols, not a description of mathematical objects.

But both views face a similar question:

\[
\text{Why does the instrument work?}
\]

\subsection*{Teaching Point}

Formalism may avoid ontology, but it inherits a serious explanatory burden: why is the
formal instrument so effective?

\subsection*{Whiteboard}

\[
\text{Instrumentalism:}
\quad
\text{science as prediction-tool}
\]

\[
\text{Formalism:}
\quad
\text{mathematics as symbol-tool}
\]

\[
\text{Question: why does it work?}
\]

\section*{10. Deductivism}

Deductivism is a more sophisticated formalist view.

\[
\boxed{
\text{Mathematics studies what follows from axioms.}
}
\]

A theorem \(\Phi\) is not simply asserted as true. Rather, the claim is:

\[
\text{If the axioms are true, then } \Phi \text{ is true.}
\]

This is why deductivism is sometimes called:

\[
\text{if-then-ism.}
\]

Deductivism preserves the importance of logical consequence. Rules of inference are not
arbitrary moves; they must preserve truth under interpretation.

\subsection*{Teaching Point}

Deductivism treats mathematics as the study of logical consequence from uninterpreted
or variably interpreted axioms.

\subsection*{Whiteboard}

\[
\textbf{Deductivism}
\]

\[
A_1,\ldots,A_n \vdash \Phi
\]

\[
\text{means: if } A_1,\ldots,A_n \text{ hold, then } \Phi \text{ holds.}
\]

\section*{11. Logical and Non-Logical Terms}

Deductivism distinguishes logical terms from mathematical primitives.

\[
\begin{array}{c|c}
\textbf{Logical terms} & \textbf{Mathematical terms} \\
\hline
\text{and} & \text{point} \\
\text{or} & \text{line} \\
\text{not} & \text{number} \\
\text{if...then} & \text{set} \\
\text{for all} & \text{function} \\
\text{there exists} & \text{plane}
\end{array}
\]

The logical vocabulary keeps its meaning. The mathematical vocabulary may be left
uninterpreted.

\subsection*{Teaching Point}

For the deductivist, mathematics does not require fixed meanings for its primitive
mathematical terms.

\subsection*{Whiteboard}

\[
\text{Keep logic fixed.}
\]

\[
\text{Let mathematical primitives vary.}
\]

\[
\text{Study consequence.}
\]

\section*{12. Hilbert's Geometry}

Hilbert's \emph{Grundlagen der Geometrie} is the central example of deductivist
mathematics.

Geometry can be axiomatized without relying on spatial intuition.

Once the axioms are stated, proofs should depend only on formal relations among the
terms.

The famous idea:

\[
\boxed{
\text{``points,'' ``lines,'' and ``planes'' could be ``tables,'' ``chairs,'' and ``beer mugs.''}
}
\]

What matters is not the intuitive meaning of the words. What matters is whether the
objects satisfy the axioms.

\subsection*{Teaching Point}

Hilbert transforms geometry from the study of intuitive space into the study of abstract
relational structure.

\subsection*{Whiteboard}

\[
\text{point, line, plane}
\]

\[
\Downarrow
\]

\[
\text{table, chair, beer mug}
\]

\[
\text{Same theory if the axioms are satisfied.}
\]

\section*{13. Implicit Definition}

Hilbert's axioms do not define individual terms one at a time.

Instead, they define the primitive terms together by fixing their relations.

This is called implicit definition.

\[
\boxed{
\text{Axioms implicitly define a structure.}
}
\]

The primitive terms get their mathematical role from the axioms.

\subsection*{Teaching Point}

Axioms can function as definitions, not by giving synonyms, but by fixing structural roles.

\subsection*{Whiteboard}

\[
\text{Explicit definition:}
\quad
\text{define one term by already understood terms.}
\]

\[
\text{Implicit definition:}
\quad
\text{fix several terms by their relations.}
\]

\[
\text{Axioms define roles.}
\]

\section*{14. Frege versus Hilbert}

Frege objects to Hilbert's use of axioms as definitions.

For Frege:

\[
\text{definitions give meanings}
\]

\[
\text{axioms state truths}
\]

So axioms cannot also be definitions.

Frege's dilemma:

\[
\begin{array}{c|c}
\textbf{Terms already have meaning} & \textbf{Terms lack meaning} \\
\hline
\text{axioms can be true or false} & \text{axioms cannot be true or false} \\
\text{but are not definitions} & \text{so they are not genuine axioms}
\end{array}
\]

Hilbert replies that this misunderstands modern axiomatics. A mathematical theory fixes
primitive notions through their relations to each other.

\subsection*{Teaching Point}

Frege wants meanings first and axioms second. Hilbert lets axioms fix the roles of the
terms.

\subsection*{Whiteboard}

\[
\begin{array}{c|c}
\textbf{Frege} & \textbf{Hilbert} \\
\hline
\text{meaning before axioms} & \text{axioms fix meaning/role} \\
\text{truth before consistency} & \text{consistency as legitimacy} \\
\text{concepts determine objects} & \text{structures determine positions}
\end{array}
\]

\section*{15. Consistency as Legitimacy}

Hilbert famously claims:

\[
\boxed{
\text{If the axioms are consistent, then the things defined by them exist.}
}
\]

A more cautious formulation is:

\[
\boxed{
\text{Consistency is enough for mathematical legitimacy.}
}
\]

This is a crucial formalist idea. Mathematics need not first identify an independently
given subject matter. It can set up a consistent axiomatic system and study its
consequences.

\subsection*{Teaching Point}

Hilbert replaces the question ``What are these objects?'' with the question ``Are the
axioms consistent?''

\subsection*{Whiteboard}

\[
\text{Existence?}
\quad \longrightarrow \quad
\text{Consistency}
\]

\[
\operatorname{Con}(T)
\quad \Rightarrow \quad
\text{legitimate theory }T
\]

\section*{16. Meta-Mathematics}

Once formal systems are rigorously specified, they can be studied mathematically.

This is meta-mathematics.

\[
\boxed{
\text{Meta-mathematics} = \text{mathematics about formal mathematics}
}
\]

Meta-mathematical questions include:

\[
\text{Is this axiom system consistent?}
\]

\[
\text{Is this axiom independent of the others?}
\]

\[
\text{Is this formula derivable from these axioms?}
\]

Hilbert's geometry already uses meta-mathematical methods by constructing models to
show consistency and independence.

\subsection*{Teaching Point}

Formalism helped create meta-mathematics: the mathematical study of formal systems
themselves.

\subsection*{Whiteboard}

\[
\text{Mathematics}
\quad \longrightarrow \quad
\text{formal system}
\]

\[
\text{Formal system}
\quad \longrightarrow \quad
\text{object of study}
\]

\[
\text{meta-mathematics}
\]

\section*{17. The Foundational Crisis}

At the turn of the twentieth century, mathematics faced both enormous progress and
serious foundational anxiety.

Progress:

\[
\text{real analysis}
\]

\[
\text{complex analysis}
\]

\[
\text{Cantor's set theory}
\]

\[
\text{the mathematics of the infinite}
\]

Crisis:

\[
\text{Russell's paradox}
\]

\[
\text{set-theoretic antinomies}
\]

\[
\text{doubts about infinite totalities}
\]

\[
\text{restrictions proposed by critics}
\]

Hilbert wanted to preserve classical mathematics without surrendering to skepticism.

\subsection*{Teaching Point}

Hilbert's mature program is a defense of classical mathematics against paradox and
restriction.

\subsection*{Whiteboard}

\[
\text{Power of the infinite}
\quad +
\quad
\text{paradoxes}
\quad =
\quad
\text{foundational crisis}
\]

\[
\text{Hilbert: defend classical mathematics.}
\]

\section*{18. Hilbert's Program}

Hilbert's goal:

\[
\boxed{
\text{Establish once and for all the certainty of mathematical methods.}
}
\]

His famous declaration:

\[
\boxed{
\text{No one shall drive us out of Cantor's paradise.}
}
\]

The strategy is to formalize ideal mathematics and then prove, using secure finitary
methods, that these formal systems are consistent.

\[
\text{formalize}
\quad \longrightarrow \quad
\text{prove consistency}
\quad \longrightarrow \quad
\text{secure ideal mathematics}
\]

\subsection*{Teaching Point}

Hilbert does not reject modern mathematics. He tries to justify it by proving that it is
safe.

\subsection*{Whiteboard}

\[
\text{Classical mathematics}
\]

\[
\Downarrow
\]

\[
\text{formal system }T
\]

\[
\Downarrow
\]

\[
\text{finitary proof of } \operatorname{Con}(T)
\]

\section*{19. Finitary Mathematics}

Hilbert distinguishes finitary mathematics from ideal mathematics.

Finitary mathematics is meaningful, secure, and directly surveyable.

It includes concrete calculations:

\[
2+3=5
\]

\[
12553+2477=15030
\]

It also includes bounded quantifiers:

\[
\exists p<101!+2 \text{ such that } p \text{ is prime.}
\]

This is finitary because only finitely many cases are involved.

\subsection*{Teaching Point}

Finitary mathematics is Hilbert's secure base: meaningful mathematics concerning
concrete symbolic constructions.

\subsection*{Whiteboard}

\[
\textbf{Finitary}
\]

\[
\text{specific calculations}
\]

\[
\text{bounded search}
\]

\[
\text{directly checkable in principle}
\]

\section*{20. Bounded and Unbounded Quantifiers}

Hilbert treats bounded and unbounded quantification differently.

\[
\begin{array}{c|c}
\textbf{Bounded quantifier} & \textbf{Unbounded quantifier} \\
\hline
\text{finite search} & \text{infinite range} \\
\text{finitary} & \text{ideal or non-finitary} \\
\text{effectively checkable in principle} & \text{not reducible to finite checking}
\end{array}
\]

Example of a bounded quantifier:

\[
\exists p<101!+2 \text{ such that } p \text{ is prime.}
\]

Example of an unbounded quantifier:

\[
\exists p>100 \text{ such that } p \text{ and } p+2 \text{ are prime.}
\]

The second statement ranges over infinitely many possible numbers.

\subsection*{Teaching Point}

The finite/infinite distinction becomes epistemologically central for Hilbert.

\subsection*{Whiteboard}

\[
\exists p<N
\quad
\text{bounded}
\]

\[
\exists p
\quad
\text{unbounded}
\]

\[
\text{finite check vs. infinite range}
\]

\section*{21. Stroke Symbols}

Hilbert identifies the natural numbers with stroke-symbols.

\[
|,\quad ||,\quad |||,\quad ||||,\ldots
\]

The symbol \(2\) abbreviates:

\[
||
\]

The symbol \(3\) abbreviates:

\[
|||
\]

The statement:

\[
3>2
\]

means that the stroke-symbol for \(3\) is longer than the stroke-symbol for \(2\).

These are not ordinary ink marks. They are symbol-types or directly surveyable symbolic
forms.

\subsection*{Teaching Point}

Hilbert's finitism has a Kantian element: it rests on immediately surveyable symbolic
constructions.

\subsection*{Whiteboard}

\[
1 = |
\]

\[
2 = ||
\]

\[
3 = |||
\]

\[
3>2
\quad
\text{because}
\quad
||| \text{ extends } ||
\]

\section*{22. Ideal Mathematics}

Everything beyond finitary mathematics is ideal mathematics.

Examples:

\[
\text{unbounded number theory}
\]

\[
\text{real analysis}
\]

\[
\text{complex analysis}
\]

\[
\text{functional analysis}
\]

\[
\text{set theory}
\]

Ideal mathematics may not have direct meaning in the way finitary mathematics does.
Rather, it is treated instrumentally, as a formal device that helps us reason efficiently.

Hilbert compares ideal mathematics to ideal points at infinity in geometry.

\subsection*{Teaching Point}

Ideal mathematics is not rejected. It is justified as a useful formal extension of finitary
mathematics.

\subsection*{Whiteboard}

\[
\begin{array}{c|c}
\textbf{Finitary Mathematics} & \textbf{Ideal Mathematics} \\
\hline
\text{meaningful} & \text{formal/instrumental} \\
\text{secure} & \text{requires justification} \\
\text{finite symbols} & \text{infinity, analysis, set theory}
\end{array}
\]

\section*{23. Conservativeness and Consistency}

Ideal mathematics is legitimate only if it does not lead to false finitary results.

This is the idea of conservativeness.

\[
\boxed{
\text{Ideal mathematics should be conservative over finitary mathematics.}
}
\]

In practice, this becomes the demand for consistency.

\[
\boxed{
\text{Prove } \operatorname{Con}(T) \text{ by finitary means.}
}
\]

If an ideal theory \(T\) is consistent, then it cannot derive contradiction. If it is also
conservative over finitary mathematics, then it cannot corrupt the meaningful base.

\subsection*{Teaching Point}

Hilbert's formalism justifies ideal mathematics by proving that it is safe for finitary
mathematics.

\subsection*{Whiteboard}

\[
\text{Ideal theory }T
\]

\[
\text{safe if}
\]

\[
T \nvdash 0 \neq 0
\]

\[
\operatorname{Con}(T)
\]

\section*{24. Von Neumann's Summary}

Von Neumann summarizes the Hilbert program in four stages:

\begin{enumerate}[leftmargin=2em]
    \item Enumerate all symbols used in mathematics and logic.
    \item Characterize the meaningful formulas.
    \item Give a construction procedure for all provable formulas.
    \item Prove, by finitary means, that the system proves the correct finitary statements.
\end{enumerate}

The first three stages concern formalization.

The fourth is the crucial epistemic stage: proving that the formalized theory is safe.

\subsection*{Teaching Point}

Hilbert's program is not merely formalization. It is formalization plus finitary justification.

\subsection*{Whiteboard}

\[
\text{symbols}
\rightarrow
\text{formulas}
\rightarrow
\text{proofs}
\rightarrow
\text{finitary consistency/safety proof}
\]

\section*{25. G\"odel's First Incompleteness Theorem}

G\"odel's first incompleteness theorem strikes at the hope for complete formalization.

Let \(T\) be a sufficiently strong, consistent, effective formal system.

Then there is a sentence \(G\) such that:

\[
T \nvdash G
\]

and, under suitable consistency assumptions,

\[
T \nvdash \neg G.
\]

Roughly, \(G\) says:

\[
\text{``}G\text{ is not provable in }T\text{.''}
\]

If \(T\) is consistent, then \(G\) is true but not provable in \(T\).

\subsection*{Teaching Point}

No sufficiently strong consistent effective formal system captures all arithmetic truth.

\subsection*{Whiteboard}

\[
T \text{ consistent}
\]

\[
\Rightarrow
\]

\[
\exists G(T\nvdash G \text{ and } T\nvdash \neg G)
\]

\[
\text{Truth outruns formal proof.}
\]

\section*{26. G\"odel's Second Incompleteness Theorem}

G\"odel's second incompleteness theorem is more directly damaging to the Hilbert
program.

No sufficiently strong consistent system can prove its own consistency.

\[
\boxed{
T \nvdash \operatorname{Con}(T)
}
\]

assuming \(T\) is consistent and sufficiently strong.

Hilbert wanted finitary consistency proofs for ideal mathematical theories.

But if the ideal theory contains enough arithmetic, then its consistency cannot be proved
inside the theory itself, much less inside a weaker finitary subsystem.

\subsection*{Teaching Point}

The second incompleteness theorem blocks the original Hilbert program's hoped-for
method of securing classical mathematics.

\subsection*{Whiteboard}

\[
\text{Hilbert wants:}
\quad
\text{finitary proof of } \operatorname{Con}(T)
\]

\[
\text{G\"odel shows:}
\quad
T \nvdash \operatorname{Con}(T)
\]

\[
\text{if }T\text{ is consistent.}
\]

\section*{27. What G\"odel Does and Does Not Show}

G\"odel's theorems do not show that mathematics is inconsistent.

They do not show that formal systems are useless.

They do not show that formalization is unimportant.

But they do show that Hilbert's original epistemic ambition cannot be fulfilled in the
simple way he hoped.

\[
\boxed{
\text{The Hilbert program cannot succeed in its original form.}
}
\]

\subsection*{Teaching Point}

G\"odel limits formalism; he does not abolish formal methods.

\subsection*{Whiteboard}

\[
\text{Not: mathematics collapses.}
\]

\[
\text{Not: formal systems are useless.}
\]

\[
\text{But: formalism cannot secure everything from within.}
\]

\section*{28. Post-G\"odel Options}

The chapter notes two possible responses after G\"odel.

First, one might challenge the formalization of consistency. There are different ways to
express consistency, and some do not fall under the exact second incompleteness result.

Second, one might argue that finitary reasoning is inherently informal and cannot be fully
captured by any one formal system.

\[
\text{Finitary reasoning}
\quad \neq
\quad
\text{one formal theory}
\]

These responses do not restore Hilbert's program in its original form, but they show that
the philosophical situation is subtle.

\subsection*{Teaching Point}

The failure of the original Hilbert program leaves room for weaker or revised formalist
projects.

\subsection*{Whiteboard}

\[
\text{Option 1: different consistency notions}
\]

\[
\text{Option 2: finitary reasoning is informal}
\]

\[
\text{Revised Hilbert programs?}
\]

\section*{29. Curry's Formalism}

The chapter ends with Haskell Curry.

Curry's central claim:

\[
\boxed{
\text{Mathematics is the science of formal systems.}
}
\]

This is different from crude game formalism. Mathematical assertions are not merely
moves within a system. They are often statements about formal systems.

A typical mathematical assertion becomes:

\[
\Phi \text{ is a theorem of } T.
\]

Curry therefore treats mathematics as objective, but its subject matter is formal systems
themselves.

\subsection*{Teaching Point}

Curry turns formalism into meta-mathematics: mathematics studies formal systems.

\subsection*{Whiteboard}

\[
\text{Mathematics}
=
\text{science of formal systems}
\]

\[
\Phi \in \operatorname{Th}(T)
\]

\[
\Phi \text{ is a theorem of }T
\]

\section*{30. Curry on Consistency and Usefulness}

Curry does not follow Hilbert in requiring consistency proofs before accepting a formal
system.

For Curry, acceptability is pragmatic.

Criteria include:

\[
\begin{array}{c}
\text{intuitive evidence of premises} \\
\text{consistency} \\
\text{usefulness}
\end{array}
\]

A consistency proof is useful, but it is not always necessary for accepting a system. If an
inconsistency appears, the system may be modified rather than abandoned completely.

\subsection*{Teaching Point}

Curry's formalism is less vulnerable to G\"odel because it does not require Hilbert-style
finitary consistency proofs.

\subsection*{Whiteboard}

\[
\text{Acceptability}
=
\text{purpose-relative}
\]

\[
\text{consistency}
+
\text{usefulness}
+
\text{evidence}
\]

\section*{31. The Problem of Informal Mathematics}

Curry's view faces a problem.

If mathematics is the science of formal systems, what about mathematics before
formalization?

Were the following people not doing mathematics?

\[
\text{Archimedes}
\]

\[
\text{Fermat}
\]

\[
\text{Euler}
\]

\[
\text{Cauchy}
\]

Even today, much mathematical practice is informal. Mathematicians do not usually write
their work in fully formal deductive systems.

\subsection*{Teaching Point}

Formalism risks mistaking the formal reconstruction of mathematics for mathematics
itself.

\subsection*{Whiteboard}

\[
\text{Formalized mathematics}
\quad \neq
\quad
\text{all mathematical practice}
\]

\[
\text{Does formalism lose the living practice of mathematics?}
\]

\section*{Conclusion: The Formalist Family}

Formalism is not one view. It is a family of views united by emphasis on symbols, rules,
axioms, formal systems, and derivations.

\[
\begin{array}{c|c|c|c}
\textbf{Term Formalism} & \textbf{Game Formalism} & \textbf{Deductivism} & \textbf{Hilbert Program} \\
\hline
\text{math about symbols} &
\text{math as symbol game} &
\text{math as consequence} &
\text{ideal math justified by consistency} \\
\text{objects = terms} &
\text{no subject matter} &
\text{if axioms, then theorem} &
\text{finitary proof of safety} \\
\text{Frege: equations fail} &
\text{Frege: application fails} &
\text{meaning thinned out} &
\text{G\"odel problem}
\end{array}
\]

Curry adds a later version:

\[
\boxed{
\text{Mathematics is the science of formal systems.}
}
\]

\section*{Final Framing}

To teach this chapter well, present formalism as both powerful and limited.

It is powerful because it clarifies proof, rigor, axioms, formal systems, consistency, and
meta-mathematics.

It is limited because mathematics appears to involve meaning, application, informal
understanding, and truth that cannot be reduced to formal derivability.

\[
\boxed{
\text{Formalism reveals the power of symbols, but struggles with meaning and application.}
}
\]

\[
\boxed{
\text{G\"odel shows that formal systems cannot fully contain arithmetic truth.}
}
\]

\section*{Discussion Questions}

\begin{enumerate}[leftmargin=2em]
    \item What motivates formalism as a philosophy of mathematics?
    \item Is mathematics really just symbol manipulation?
    \item What is the difference between term formalism and game formalism?
    \item Can numbers be identified with numerals?
    \item Why does Frege think formalism cannot explain applicability?
    \item Is mathematics more like chess or more like science?
    \item What does deductivism add to simple game formalism?
    \item Can axioms function as definitions?
    \item Was Frege or Hilbert right about the relation between definitions and axioms?
    \item Is consistency enough for mathematical existence?
    \item What is the difference between finitary and ideal mathematics?
    \item Why did Hilbert want finitary consistency proofs?
    \item What does G\"odel's first incompleteness theorem show?
    \item Why is the second incompleteness theorem a problem for Hilbert's program?
    \item Is Curry right that mathematics is the science of formal systems?
\end{enumerate}

\section*{Key Terms}

\begin{description}[leftmargin=2em]
    \item[Formalism] A family of views emphasizing mathematics as symbol manipulation,
    formal rules, axioms, and formal systems.

    \item[Term formalism] The view that mathematical objects are symbols or terms.

    \item[Game formalism] The view that mathematics is like a game played with symbols
    according to rules.

    \item[Token] A particular physical inscription or mark.

    \item[Type] A repeatable abstract symbol-form shared by many tokens.

    \item[Deductivism] The view that mathematics studies what follows logically from
    axioms.

    \item[If-then-ism] Another name for deductivism; the view that mathematical theorems
    are conditional on axioms.

    \item[Implicit definition] A definition in which axioms fix the roles of primitive terms
    by specifying their relations.

    \item[Consistency] The property of a formal system that it does not derive a
    contradiction.

    \item[Meta-mathematics] The mathematical study of formal mathematical systems.

    \item[Finitary mathematics] For Hilbert, secure mathematics concerning directly
    surveyable finite symbolic constructions.

    \item[Ideal mathematics] Mathematics involving ideal elements such as infinity, set
    theory, and analysis, justified instrumentally by consistency or conservativeness.

    \item[Conservativeness] The property that an extension of a theory does not prove any
    new claims of a specified base kind, especially no false finitary claims.

    \item[Hilbert program] Hilbert's project of formalizing ideal mathematics and proving
    its consistency by finitary means.

    \item[G\"odel's first incompleteness theorem] The result that any sufficiently strong,
    consistent, effective formal system is incomplete.

    \item[G\"odel's second incompleteness theorem] The result that no sufficiently strong
    consistent formal system can prove its own consistency.

    \item[Curry's formalism] The view that mathematics is the science of formal systems.
\end{description}

\end{document}